% Abstract. Nature summary-paragraph move order (nature-summary-paragraph.pdf):
% basic introduction / detailed background / general problem / main result ("here we
% show" equivalent) / what it reveals / general context / broader perspective.
% Written to the submission portal's 3550-character limit rather than Nature's 300-word
% one, so each move gets a paragraph instead of running them together.
% Impersonal voice throughout. Assumes a computer-engineering reader with no specialism
% in evolutionary computation: the method is named and the jargon glossed.
% Does not characterise other people's systems. The novelty is put as a claim about this
% work ("not a new idea, though it has not been taken far"), not as a comparison.
% Every score carries the budget that produced it.
% Must match Chapter 6 and content/summary.tex.
% Regenerate the portal copy with: python3 tools/make-abstract-plaintext.py

Machines write working code and pass professional exams, yet fail at puzzles most people
solve without training. The Abstraction and Reasoning Corpus (ARC-AGI) makes that gap
measurable. Each task gives a few coloured grids and asks for the rule between them, drawing
only on notions that developmental psychology considers innate: objects, geometry, counting,
and space.

Because every task is new, no training distribution can be fitted and data contamination
cannot inflate a score. Humans average above
seventy percent, while plain language models manage about ten; reasoning
models reach above ninety, but on subsets of the tasks and at inference costs of dollars per
puzzle. The symbolic route is to search for a short program fitting the examples, which
makes the difficulty one of representation, since over raw pixels the programs needed are
hopelessly long.

The question this thesis takes up is whether that representation can itself be a large
parameterized space, searched per task alongside the program and within a budget that fits one
machine.

This thesis shows that it can. An evolutionary algorithm co-evolves three components per
task: a \emph{parser} turning a grid into a typed graph of objects and relations, a short
\emph{program} of rewrite rules over that graph, and a \emph{projector} rendering the result
back. The current parser is seventeen detection passes under forty-four evolvable parameters, in the same
genome as the program. Candidates are ranked first on reproducing every train pair
exactly, then by a graph distance over those that fail and a leave-one-out penalty favouring
programs that survive dropping any pair. The best single run, in the solver's C
implementation, fits the train pairs of 211 tasks in the ARC-AGI-1 training set and reproduces
the withheld pair on 163 of the 391 it attempted, the count rising with the compute given. It
finishes on one ordinary machine, where the models scoring higher on this benchmark run in
data centres at dollars per puzzle. Pooling runs of differing engines, budgets, and seeds, 217
tasks fall to at least one run at two attempts, though no single run reaches that.

The winning programs are tiny: one rewrite rule at the median, nine in ten using at most two,
which is the point of co-evolving the representation, since the parser absorbs modelling a
pixel-level language would spell out. Two findings then cut against ``expectation''. The limit
is not perception, since tasks narrowly missed are
parsed more consistently than those solved; it is the expressiveness of the rules. And the gap between that 217 and any one run is a
failure of choice, not of search: fitting programs are found reliably, but the ranking cannot
always tell which will generalize.

The ceiling is what the rules can express, not the machine. On ARC-AGI-2, built to demand the
multi-step composition a one-rule program cannot express, the solver fits 3 of the 120 public
evaluation tasks and solves 2, one of those already in ARC-AGI-1. Mining reusable strategies,
which do recur across a third of held-out tasks, is one attempt at that gap and does not yet
pay.

Solving ARC-AGI is not solving artificial general intelligence, and the interest lies less in
the score than in the route. Each solution is a short program that can be read and
checked rather than a set of weights, and the run fits on ordinary hardware, so it suits
settings where a decision must be inspected or compute is scarce.
Interactive environments are the next step, with the same graph predicted forward instead of
rewritten once.
